%!%
% Copyright (c) 2026 mingcheng <mingcheng@outlook.com>
% 
% File: ch11-troubleshooting.tex
% Author: mingcheng <mingcheng@outlook.com>
% File Created: 2026-03-27 14:57:14
% 
% Modified By: mingcheng <mingcheng@outlook.com>
% Last Modified: 2026-03-29 07:50:49
%%

% !TeX root = ../prompt-engineering-guide-for-beginners.tex

\chapter{翻车救援指南：模型不听话怎么办？}
\label{ch:troubleshooting}

在驱使 AI 当牛做马的过程中，你大概率会遇上各种各样的「翻车名场面」：回答得牛头不对马嘴、大张旗鼓地捏造数据、产出的文案像裹脚布一样又臭又长且毫无灵魂……
请安抚你的情绪，这些病症几乎都有特效解药。本章将手把手带你建立急救应对包，帮你逐一排除雷区。

\section{AI 开始一本正经地胡说八道（幻觉）怎么办？}

大模型最常见的致命劣根性被称为「幻觉」（Hallucination）。它的表现症状是：用极为扎实且自信的语气，堂而皇之地向你编造完全不存在的人名、事件、学术名言乃至莫须有的财务数据。

为什么会这样？因为 AI 的本质是「生成看起来合理的文字」，而不是「查找准确的事实」。当它不知道答案时，它不会说「我不知道」，而是会「编一个」。

\textbf{应对策略：}

\begin{itemize}
  \item \textbf{在提示词中设定边界}：「请仅根据我提供的文档内容来回答。如果文档中没有相关信息，请直接说『我无法确认』，不要猜测或补充。」
  \item \textbf{用文档替代记忆}：把你手头的权威资料、操作手册或内部文件直接粘贴到对话中，让 AI 基于这些有据可查的内容来回答。
  \item \textbf{要求标注信息来源}：「请在你的结论中标注每个观点的依据来源（来自哪个机构、哪一年的数据等）。」这虽然不能完全杜绝捏造，但能帮你快速核实真伪。
  \item \textbf{重要信息务必人工核实}：涉及公司宣传、法律条款、财务数据等零容错的场景，一定要用搜索引擎或其他可靠来源进行交叉验证，不要直接采信 AI 的输出。
\end{itemize}

\section{回答太长、太空洞、太像「机器人」怎么办？}

有时候 AI 为了显得全面，会输出一大段字数不少但缺乏实际价值的内容，语言里充斥着「首先、其次、总而言之、希望对你有所帮助」等套话，读起来像是在做填充题。

\textbf{解决办法：}

\begin{itemize}
  \item \textbf{设定字数限制}：「回答控制在 100 字以内」「用不超过三个要点来回答」。
  \item \textbf{指定你想要的文风}：「用知乎高赞回答的风格来写，简洁有力，带点实际经验」。
  \item \textbf{禁止套话}：「不要在开头和结尾加任何客套话或总结性废话，直接给出核心内容」。
  \item \textbf{要求可执行的干货}：「跳过背景介绍，直接给我三条明天就能落地执行的具体建议」。
\end{itemize}

\begin{tipbox}
常备的救场句式：「你这版回答太冗长了，帮我浓缩成核心观点，去掉开场白和结尾的客套话。」
\end{tipbox}

来看一个具体的对比。假设你想让 AI 给你写一封拒绝加班的回复，如果不加任何约束：

\begin{promptbox}
帮我写一条拒绝加班的消息发给领导。
\end{promptbox}

AI 可能会给你一大段客客气气、反复表决心的冗长文字。现在加上约束：

\begin{promptbox}
帮我写一条拒绝加班的消息发给领导。要求：不超过 50 字，语气礼貌但坚定，不要解释原因，直接表明今晚有事无法加班。
\end{promptbox}

\begin{aibox}
王总您好，今晚已有无法调整的安排，加班恐怕难以配合，抱歉给您添麻烦。如明日有需要我可以优先处理。
\end{aibox}

同样的需求，加上字数、语气和内容的约束后，输出就从「裹脚布」变成了可以直接发出去的消息。

\section{挤牙膏式追问法}

AI 第一次的回答不够好是很正常的。好消息是，你可以通过多轮对话不断追问和微调，逐步「逼」出你想要的结果。

追问的关键在于：告诉 AI 哪里不满意，以及你想要什么样的改变。不要只说「不好」「重新写」，而是说清楚「哪里不好」「怎样才算好」。想象一下你在指导一个实习生修改文档的场景，你不会只说「这个不行」，而会说「第二段的语气太生硬了，改成更口语化的表达」。对 AI 也是一样的道理。

\textbf{常用的追问句式：}

\begin{itemize}
  \item 「这个方向不错，但请在第二点上再展开一些，加入具体案例。」
  \item 「语气太正式了，请改成更轻松的风格。」
  \item 「你漏掉了XX方面，请补充进去。」
  \item 「这个回答太理论化了，请用一个实际工作中的例子来说明。」
  \item 「请把上面的内容重新组织一下，按照重要性从高到低排列。」
  \item 「不够具体。请给出可以直接执行的步骤，而不是泛泛的建议。」
\end{itemize}

这些追问并不需要从头重写提示词，只是在现有对话的基础上继续调整。通常经过二到三轮追问，你就能拿到相当满意的结果了。如果追问了五六轮还是不行，那往往说明你的初始提示词需要重新设计，而不是继续「挤牙膏」。

\section{Prompt 迭代优化流程}

写提示词不是一次性的事情，而是一个不断迭代优化的过程。就像写代码需要调试（Debug）一样，好的提示词也是改出来的：

\begin{enumerate}
  \item \textbf{起草（Draft）}：先把核心需求和原始材料直接写进提示词中，不必纠结措辞，先发出去试水。
  \item \textbf{观察和诊断（Inspect）}：看 AI 返回的结果，找出问题出在哪里。是理解有偏差？段落太散？还是风格不对？
  \item \textbf{针对性修改（Iterate）}：回到提示词中，利用前面章节学到的技巧，针对具体问题追加约束或补充说明，然后让它重新生成。
  \item \textbf{验证和打磨（Verify）}：如果还有瑕疵，继续「写 $\to$ 看效果 $\to$ 改」这个循环，直到满意为止。
\end{enumerate}

\begin{tipbox}
极其高亮的一条实验铁律：\textbf{你如果在进行调优，请保证绝不贪刀，尽量每次仅为提示词修正并验证「一个维度的变量因素」}（不要一拍脑门同时把语气、字数和格式全加上）。当你同时混杂着做四五个颠覆性的指令改动时，万一产出品质全面雪崩暴走，你作为控制者根本无从排查是被哪一个草率的参数调整拖垮了整盘棋局。
\end{tipbox}

另外一个好习惯是：\textbf{为自己的顺手武器库打版}。当你在经过七八轮焦灼的推拉后，终于在这个极其复杂的场景榨取到了想要的那条满状态金牌提示词大纲，别忘了把它放进你的便签本供起来（就像程序员积攒极佳的脚手架代码那样），下次直接调用。这也是后续章节「搭建私有指令库」的底层逻辑。

\section{AI 的回答自相矛盾？}

你可能会发现：同样的问题，第一次问和第二次问的答案不一样。这不是 Bug，而是 AI 的正常行为。

原因在于第~\ref{ch:basic-concepts}~章提到的 Temperature 参数（参见第~\ref{sec:temperature}~节）。AI 的回答本质上是基于概率的「采样」，每次采样的结果可能不同，就像掷骰子一样。

\textbf{应对方法：}

\begin{itemize}
  \item 如果你需要稳定一致的回答（比如数据分析、分类任务），在支持参数设置的工具中把 Temperature 调低
  \item 用「自我一致性」技巧（参见第~\ref{sec:self-consistency}~节），让 AI 回答多次，取最一致的结果
  \item 在提示词中明确要求：「请给出最确定的答案，不要给出模棱两可的回复」
  \item 对于重要决策，多问几次，看看答案是否一致。如果不一致，说明这个问题可能超出了 AI 能可靠回答的范围
\end{itemize}

\section{长文本被「截断」或「遗忘」了？}

当你输入很长的文本让 AI 处理，或者和 AI 聊了很多轮之后，你可能发现它开始「忘记」之前说过的内容。

这是因为上下文窗口有大小限制（参见第~\ref{sec:context-window}~节）。当总内容超过上限时，最早的内容就会被丢弃。

\textbf{应对方法：}

\begin{itemize}
  \item \textbf{分段处理长文本}：不要一次性把整本书贴进去。把长文本分成几个段落，每次处理一段，最后让 AI 汇总。
  \item \textbf{先总结再深入}：如果需要 AI 分析一篇很长的文章，可以先让它写一个总结，然后基于总结再深入讨论。
  \item \textbf{适时开新对话}：如果当前的话题已经结束，开一个新的对话窗口，把上下文空间腾出来。
  \item \textbf{重复关键信息}：如果你在对话进行中发现 AI 「忘了」某些重要设定，直接再说一遍就行。比如「请记住，你的角色是一位严格的编辑，请继续用这个身份回答我的问题。」
  \item \textbf{使用支持长上下文的模型}：如果你经常需要处理长文本，选择上下文窗口更大的模型（比如 Kimi 支持超长文本，Claude 也支持很长的上下文）。
\end{itemize}

总结一下这章的核心心法：AI 不是用一次就完美的工具，它更像是一个需要你不断调教和引导的合作伙伴。遇到问题不要放弃，试着调整你的提示词，大多数「翻车」都能救回来。

解决了「怎么用好」的问题，接下来我们还需要聊一个同样重要的话题：怎么「用得安全」。下一章会讨论使用 AI 时的隐私保护、事实核查和伦理边界。
